% TOP_PROBLEM: {"id": 6200026, "title": "Boundaries of Groups and Kleinian Groups — Problem 26", "category_id": 6, "category": {"id": 6, "name": "geometry", "display_name": "Geometry", "description": "Euclidean and non-Euclidean geometry, geometric structures.", "slug": "geometry", "order_index": 6, "created_at": "2026-07-31T15:26:25.671Z"}, "status": "open", "problem_number": "AMR-061-0026", "set_id": 13}
% FIRST_POSTED: 2026-10-02
% ENTRY_KIND: statement_only
% UnsolvedMath ID 6200026; source catalogue code AMR-061-0026
% !TeX program = lualatex
% Definition and sources only; this file is not an LLM solution attempt.
\documentclass[11pt,a4paper]{article}
\usepackage[margin=25mm]{geometry}
\usepackage{fontspec}
\setmainfont{lmroman10-regular.otf}[BoldFont=lmroman10-bold.otf,
 ItalicFont=lmroman10-italic.otf,BoldItalicFont=lmroman10-bolditalic.otf]
\usepackage{amsmath,amssymb,mathtools}
\usepackage{unicode-math}
\setmathfont{latinmodern-math.otf}
\AtBeginDocument{\let\setminus\smallsetminus}
\usepackage{xurl,hyperref}
\hypersetup{colorlinks=true,urlcolor=blue}
\setcounter{secnumdepth}{0}
\setlength{\parindent}{0pt}
\setlength{\parskip}{5pt}
\setlength{\emergencystretch}{3em}
\begin{document}
\section*{Boundaries of Groups and Kleinian Groups — Problem 26}\label{6200026}
{\small UnsolvedMath ID 6200026 · AMR-061-0026 · Geometry · AMR Open Problem Lists}
\subsection{Definitions and mathematical statement}
A $Z$-structure for a group $G$ is a compactification $\overline X=X\cup Z$ in which $\overline X$ is a compact Euclidean retract, $Z$ is a $Z$-set, and $G$ acts freely, properly and cocompactly on $X$. Being a $Z$-set means that the identity of $\overline X$ can be homotoped instantaneously into $X$. One also requires nullity: translates of each fixed compact subset of $X$ eventually become arbitrarily small with respect to every open cover of $\overline X$.

Use the explicit cell-like comparison in the source: two boundaries have a common compact metrizable domain mapping continuously and surjectively onto each, with both maps cell-like. A surjection is cell-like when every fiber has the shape of a point, meaning it contracts inside each neighborhood in an ambient absolute neighborhood retract.

Can one group $G$ have two $Z$-structure boundaries $Z_1,Z_2$ which are not cell-like equivalent? Equivalently, must every pair of boundaries admitted by these axioms have such a common domain and such maps?

The spaces giving the two structures may be different, and the boundary actions need not extend: extension is an additional $EZ$ condition, not part of this question. A homeomorphism would give cell-like equivalence, but the converse need not hold. Therefore examples of nonhomeomorphic boundaries do not settle the problem. The question concerns the stated stronger relation beyond the shape-equivalence information already available for $Z$-boundaries.
\subsection{Short English statement}
Can the same group have two Z-boundaries that remain different even up to cell-like equivalence?
\subsection{Sources}
\begin{itemize}
\item[S1] ULAM.AI and the UnsolvedMath Contributors, supplied problems.json, record 6200026 (AMR-061-0026), AMR Open Problem Lists. Catalogue identity and source transcription; CC BY 4.0.
\url{https://huggingface.co/datasets/ulamai/UnsolvedMath}.
\item[S2] Kapovich - Problems on Boundaries of Groups and Kleinian Groups (2005), Problem 26, PDF page 8. Original source indexed by AMR-061-0026.
\url{https://www.math.ucdavis.edu/~kapovich/EPR/problems.pdf}.
\end{itemize}
\end{document}
